\section{The Dynamical Principle: Transitions as Entropic Action Minima}
\label{sec:dynamical_principle}

In \textit{Informational Energetics: Entropic Action} \cite{meyer_ieaction_2026}, we formalized the universal requirement for persistence as the minimization of the Entropic Action: $S_\Phi = \int \mathcal{D} d\mu$. The instantaneous dissipation density is the product of three universal factors:
\begin{equation}
    \mathcal{D} = \xi \cdot \mathbf{K} \cdot f
\end{equation}
where $\xi$ is the transition impedance (the vacuum coupling $\alpha^{-1}$), $\mathbf{K}$ is the topological mismatch or structural stress (defined by invariants like the boundary $\chi$ and manifold $D$), and $f$ is the transition probability or rate (the mixing angle).

To persist, the system minimizes transition amplitudes ($f$) strictly to the level required to dissipate internal structural stress ($\mathbf{K}$) against the vacuum impedance ($\xi$). When applied to the $E_8$ substrate, this dynamic minimization yields three distinct transition archetypes:
\begin{enumerate}
    \item \textbf{Type I (Chiral Leak):} Adjacent linear transitions governed by minimal interface cost.
    \item \textbf{Type II (Dimensional Shift):} Volumetric promotions engaging the full field impedance.
    \item \textbf{Type III (Quantum Tunneling):} Non-adjacent transitions across a maximal geometric barrier.
\end{enumerate}

\textbf{Validation via QED:} To demonstrate that this geometric impedance logic is universal and not merely over-fitted to the flavor sector, we apply it to standard Quantum Electrodynamics. In Appendix \ref{sec:impedancepowertransfer}, we demonstrate that treating interactions as geometric power transfers naturally derives the $e^+e^- \to \mu^+\mu^-$ scattering cross-section and exactly recovers Schwinger's anomalous magnetic moment ($a_e = \alpha/2\pi$) without loop integrals.